% O011-BR — modul asli, bukan terjemahan atau adaptasi dari teks pembanding.
% Judul: Jembatan Kohomologi de Rham dan Topologi Diferensial
% Lisensi modul: CC BY-SA 4.0.
% Provenans: disusun oleh OpenAI Codex gpt-5.6-sol, Ultra, atas arahan pengguna.
% Prasyarat internal: Unit 14, 20, 22, dan 23 edisi ini.

\setcounter{section}{31}
\section{Jembatan Kohomologi de Rham dan Topologi Diferensial}
\label{o011-bridge-de-rham}

\subsection{Tujuan dan prasyarat}

Unit 14 memperkenalkan bentuk diferensial, Unit 20 turunan eksterior, Unit 22
partisi kesatuan, dan Unit 23 teorema Stokes. Modul ini menyatukan keempatnya
menjadi sebuah invariant topologis. Gagasan pokoknya sederhana: bentuk tertutup
yang bukan turunan eksterior bentuk lain mendeteksi ``lubang'' pada manifold.
Kita membuktikan hasil lokal, membangun invariansi homotopi, memakai barisan
Mayer--Vietoris untuk perhitungan global, lalu menandai jalan menuju derajat,
transversalitas, dan teori Morse tanpa berpura-pura membahas teori-teori itu
secara lengkap.

Semua manifold di bagian ini dianggap mulus, berdimensi hingga, Hausdorff, dan
memiliki basis terhitung. Ruang bentuk-$k$ mulus pada $M$ ditulis
$\Omega^k(M)$; untuk $k<0$ atau $k>\dim M$, ruang itu adalah nol.

\subsection{Kompleks de Rham}

\begin{definition}\label{o011-br-def-closed-exact}
Suatu bentuk $\omega\in\Omega^k(M)$ disebut \emph{tertutup} jika
$d\omega=0$, dan disebut \emph{eksak} jika terdapat
$\eta\in\Omega^{k-1}(M)$ dengan $\omega=d\eta$.
\end{definition}

Karena $d\circ d=0$, setiap bentuk eksak tertutup. Jadi turunan eksterior
membentuk kompleks rantai-kobalik
\[
  0\longrightarrow\Omega^0(M)\xrightarrow{d}\Omega^1(M)
  \xrightarrow{d}\cdots\xrightarrow{d}\Omega^n(M)
  \longrightarrow0.
\]

\begin{definition}\label{o011-br-def-cohomology}
Kohomologi de Rham derajat $k$ adalah ruang vektor
\[
  H^k_{\mathrm{dR}}(M)
  =\frac{\ker(d\colon\Omega^k(M)\to\Omega^{k+1}(M))}
  {\operatorname{im}(d\colon\Omega^{k-1}(M)\to\Omega^k(M))}.
\]
Kelas bentuk tertutup $\omega$ ditulis $[\omega]$.
\end{definition}

\begin{Proposition}\label{o011-br-prop-wedge-ring}
Produk wedge menurunkan produk bergradasi
\[
 H^p_{\mathrm{dR}}(M)\times H^q_{\mathrm{dR}}(M)
 \longrightarrow H^{p+q}_{\mathrm{dR}}(M),
 \qquad ([\alpha],[\beta])\longmapsto[\alpha\wedge\beta].
\]
Produk ini memenuhi
$[\alpha]\wedge[\beta]=(-1)^{pq}[\beta]\wedge[\alpha]$.
\end{Proposition}

\begin{proof}
Aturan Leibniz bergradasi memberi
\[
 d(\alpha\wedge\beta)=d\alpha\wedge\beta
 +(-1)^p\alpha\wedge d\beta.
\]
Jika $\alpha$ dan $\beta$ tertutup, wedge-nya tertutup. Jika
$\alpha$ diganti $\alpha+d\mu$, maka
\[
 (\alpha+d\mu)\wedge\beta-\alpha\wedge\beta
 =d(\mu\wedge\beta)
\]
karena $d\beta=0$. Argumen yang sama berlaku pada faktor kedua. Jadi produk
tidak bergantung pada wakil. Komutativitas bergradasi diwarisi dari bentuk.
\end{proof}

\begin{Proposition}\label{o011-br-prop-h0-components}
$H^0_{\mathrm{dR}}(M)$ adalah ruang fungsi yang konstan pada setiap komponen
terhubung. Khususnya, jika $M$ terhubung, $H^0_{\mathrm{dR}}(M)\cong\mathbb R$.
\end{Proposition}

\begin{proof}
Untuk fungsi mulus $f$, persamaan $df=0$ berarti semua turunan arah lokalnya
nol. Maka $f$ konstan pada setiap komponen terhubung. Tidak ada bentuk
berderajat $-1$, sehingga pada derajat nol tidak ada hasil bagi tambahan.
\end{proof}

\subsection{Tarikan balik dan funktorialitas}

Jika $f\colon M\to N$ mulus, tarikan balik memenuhi
\[
 f^*(\alpha\wedge\beta)=f^*\alpha\wedge f^*\beta,
 \qquad d(f^*\alpha)=f^*(d\alpha).
\]

\begin{Proposition}\label{o011-br-prop-pullback-cohomology}
Pemetaan $f$ menginduksi homomorfisme aljabar bergradasi
\[
 f^*\colon H^*_{\mathrm{dR}}(N)\longrightarrow H^*_{\mathrm{dR}}(M).
\]
Jika $g\colon N\to P$ mulus, maka $(g\circ f)^*=f^*\circ g^*$, dan
$\operatorname{id}_M^*$ adalah identitas.
\end{Proposition}

\begin{proof}
Komutasi dengan $d$ menunjukkan bahwa tarikan balik membawa bentuk tertutup
ke bentuk tertutup dan bentuk eksak ke bentuk eksak. Karena itu ia terdefinisi
pada kelas kohomologi. Semua identitas lain sudah berlaku pada tingkat bentuk.
\end{proof}

Perhatikan pembalikan arah: pemetaan $M\to N$ menghasilkan pemetaan
kohomologi dari $N$ ke $M$. Kohomologi de Rham adalah funktor kontravarian.

\subsection{Lema Poincar\'e pada himpunan berbentuk bintang}

Suatu himpunan terbuka $U\subseteq\mathbb R^n$ disebut berbentuk bintang
terhadap $0$ jika $tx\in U$ untuk semua $x\in U$ dan $t\in[0,1]$. Misalkan
$E=\sum_i x_i\partial_{x_i}$ adalah medan Euler dan $\iota_E$ kontraksi oleh
$E$.

Untuk $\omega\in\Omega^k(U)$ dengan $k\ge1$, definisikan operator homotopi oleh
\[
 (K\omega)_x(v_1,\ldots,v_{k-1})
 =\int_0^1 t^{k-1}\omega_{tx}(x,v_1,\ldots,v_{k-1})\,dt.
\]
Dengan medan Euler $E_y=y$, integran yang sama dapat ditulis
$t^{k-2}(\iota_E\omega)_{tx}$ untuk $t>0$; rumus pertama tetap teratur di
$t=0$ dan karena itu kita pakai sebagai definisi. Rumus yang lebih intrinsik
akan muncul pada homotopi umum di bawah.

\begin{Theorem}[Lema Poincar\'e]\label{o011-br-thm-poincare}
Jika $U\subseteq\mathbb R^n$ terbuka dan berbentuk bintang, setiap bentuk
tertutup $\omega\in\Omega^k(U)$ dengan $k\ge1$ adalah eksak. Dengan kata lain,
$H^k_{\mathrm{dR}}(U)=0$ untuk $k\ge1$.
\end{Theorem}

\begin{proof}
Ambil homotopi $F\colon U\times[0,1]\to U$, $F(x,t)=tx$. Perhitungan di
koordinat, atau rumus homotopi yang dibuktikan pada Teorema
\ref{o011-br-thm-chain-homotopy}, memberi
\[
 dK\omega+K(d\omega)=F_1^*\omega-F_0^*\omega.
\]
Untuk $k\ge1$, tarikan balik bentuk-$k$ oleh pemetaan konstan $F_0$ adalah nol,
sedangkan $F_1$ adalah identitas. Jika $d\omega=0$, maka
$\omega=d(K\omega)$.
\end{proof}

\begin{bemerkung}\label{o011-br-rem-local-global}
Lema Poincar\'e bersifat lokal: setiap titik manifold mempunyai lingkungan
koordinat yang dapat diperkecil menjadi bola, sehingga setiap bentuk tertutup
berderajat positif lokalnya eksak. Hambatan untuk memperoleh primitif global
adalah tepat informasi yang diukur kohomologi de Rham.
\end{bemerkung}

\subsection{Homotopi rantai dan invariansi homotopi}

Misalkan $F\colon M\times[0,1]\to N$ suatu homotopi mulus dan
$F_t(p)=F(p,t)$. Untuk $\omega\in\Omega^k(N)$, definisikan
\[
 K_F\omega
 =\int_0^1 \iota_t^*\bigl(\iota_{\partial_t}F^*\omega\bigr)\,dt
 \in\Omega^{k-1}(M),
\]
di mana $\iota_t(p)=(p,t)$.

\begin{Theorem}[Rumus homotopi rantai]\label{o011-br-thm-chain-homotopy}
Operator $K_F$ memenuhi
\[
 dK_F+K_Fd=F_1^*-F_0^*.
\]
\end{Theorem}

\begin{proof}
Pada $M\times[0,1]$, turunan terhadap $t$ dari tarikan balik sepanjang
$\iota_t$ memenuhi rumus Cartan
\[
 \frac d{dt}\iota_t^*(F^*\omega)
 =\iota_t^*\mathcal L_{\partial_t}(F^*\omega)
 =\iota_t^*\bigl(d\iota_{\partial_t}F^*\omega
 +\iota_{\partial_t}dF^*\omega\bigr).
\]
Integrasikan dari $0$ sampai $1$, gunakan bahwa tarikan balik komutatif dengan
$d$, dan kenali kedua suku sebagai $dK_F\omega$ serta $K_Fd\omega$.
\end{proof}

\begin{Korollar}\label{o011-br-cor-homotopy-invariance}
Pemetaan mulus yang homotopik menginduksi pemetaan kohomologi yang sama. Jika
$M$ dan $N$ ekuivalen-homotopi, maka
$H^*_{\mathrm{dR}}(M)\cong H^*_{\mathrm{dR}}(N)$ sebagai aljabar bergradasi.
\end{Korollar}

\begin{proof}
Jika $d\omega=0$, rumus homotopi memberi
$F_1^*\omega-F_0^*\omega=d(K_F\omega)$, sehingga kedua tarikan balik menentukan
kelas yang sama. Untuk ekuivalensi homotopi, terapkan funktorialitas pada
komposisi yang homotopik dengan identitas.
\end{proof}

\begin{Korollar}\label{o011-br-cor-contractible}
Jika $M$ kontraktibel, maka $H^0_{\mathrm{dR}}(M)\cong\mathbb R$ dan
$H^k_{\mathrm{dR}}(M)=0$ untuk $k>0$.
\end{Korollar}

\subsection{Barisan Mayer--Vietoris}

Misalkan $M=U\cup V$ dengan $U,V$ terbuka. Untuk setiap $k$ terdapat barisan
pendek kompleks
\[
0\longrightarrow\Omega^k(M)
\xrightarrow{r}(\Omega^k(U)\oplus\Omega^k(V))
\xrightarrow{\delta}\Omega^k(U\cap V)
\longrightarrow0,
\]
dengan
\[
 r(\omega)=(\omega|_U,\omega|_V),
 \qquad \delta(\alpha,\beta)=\alpha|_{U\cap V}-\beta|_{U\cap V}.
\]
Injektivitas dan ketepatan di tengah mengikuti sifat pengeleman bentuk. Untuk
surjektivitas $\delta$, pilih partisi kesatuan $\rho_U+\rho_V=1$ yang
tersubordinasi pada $U,V$; suatu bentuk $\gamma$ di $U\cap V$ diangkat oleh
perpanjangan nol yang dibangun dari $\rho_V\gamma$ pada $U$ dan
$-\rho_U\gamma$ pada $V$.

\begin{Theorem}[Mayer--Vietoris]\label{o011-br-thm-mv}
Penutup $M=U\cup V$ menghasilkan barisan eksak panjang
\[
\cdots\longrightarrow H^{k-1}_{\mathrm{dR}}(U\cap V)
\xrightarrow{\partial}H^k_{\mathrm{dR}}(M)
\xrightarrow{r^*}H^k_{\mathrm{dR}}(U)\oplus H^k_{\mathrm{dR}}(V)
\xrightarrow{\delta^*}H^k_{\mathrm{dR}}(U\cap V)
\xrightarrow{\partial}H^{k+1}_{\mathrm{dR}}(M)
\longrightarrow\cdots.
\]
\end{Theorem}

\begin{proof}
Barisan pendek kompleks di atas menghasilkan barisan eksak panjang kohomologi
melalui konstruksi standar pemetaan penghubung. Secara konkret, jika
$\gamma$ tertutup pada $U\cap V$, pilih $(\alpha,\beta)$ dengan
$\delta(\alpha,\beta)=\gamma$. Karena $d\gamma=0$, pasangan
$(d\alpha,d\beta)$ saling cocok pada irisan, sehingga menempel menjadi bentuk
tertutup $\omega$ pada $M$. Tetapkan $\partial[\gamma]=[\omega]$. Perubahan
pilihan hanya mengubah $\omega$ dengan bentuk eksak, dan pemeriksaan langsung
memberi ketepatan pada setiap suku.
\end{proof}

\subsection{Perhitungan dasar}

\begin{Proposition}\label{o011-br-prop-circle}
Kohomologi de Rham lingkaran adalah
\[
 H^0_{\mathrm{dR}}(S^1)\cong\mathbb R,
 \qquad H^1_{\mathrm{dR}}(S^1)\cong\mathbb R,
 \qquad H^k_{\mathrm{dR}}(S^1)=0\quad(k\ge2).
\]
Isomorfisme pada derajat satu dapat diberikan oleh periode
\[
 [\omega]\longmapsto\int_{S^1}\omega.
\]
\end{Proposition}

\begin{proof}
Derajat nol mengikuti keterhubungan, dan derajat di atas satu nol karena
dimensi. Tutupi $S^1$ oleh dua busur terbuka kontraktibel $U,V$ yang irisannya
mempunyai dua komponen. Bagian relevan Mayer--Vietoris adalah
\[
 0\to H^0(S^1)\to\mathbb R\oplus\mathbb R
 \to\mathbb R^2\to H^1(S^1)\to0.
\]
Pemetaan tengah mengirim $(a,b)$ ke $(a-b,a-b)$, sehingga cokernelnya satu
dimensi. Integral lenyap pada bentuk eksak dan tidak nol pada bentuk sudut
ternormalisasi; karena itu periode memberi isomorfisme yang dinyatakan.
\end{proof}

Pada $\mathbb R^2\setminus\{0\}$, bentuk
\[
 \eta=\frac{1}{2\pi}\frac{x\,dy-y\,dx}{x^2+y^2}
\]
tertutup dan integralnya pada lingkaran satuan adalah $1$. Jadi ia tidak eksak
dan mewakili generator kelas derajat satu.

\begin{Theorem}\label{o011-br-thm-sphere}
Untuk $n\ge1$,
\[
 H^k_{\mathrm{dR}}(S^n)\cong
 \begin{cases}
 \mathbb R,&k=0\text{ atau }k=n,\\
 0,&\text{selain itu}.
 \end{cases}
\]
\end{Theorem}

\begin{proof}
Kasus $n=1$ sudah dibuktikan. Untuk $n\ge2$, tutupi $S^n$ oleh lingkungan
terbuka kutub utara dan selatan yang masing-masing kontraktibel dan yang
irisannya ekuivalen-homotopi dengan $S^{n-1}$. Mayer--Vietoris, di luar
derajat nol, memberi isomorfisme penghubung
\[
 H^{k-1}_{\mathrm{dR}}(S^{n-1})\cong H^k_{\mathrm{dR}}(S^n)
 \qquad(1<k\le n).
\]
Pada derajat satu, keterhubungan penutup dan irisan memberi nol. Induksi pada
$n$ menghasilkan satu kelas di derajat teratas dan nol pada derajat antara.
Derajat nol kembali mengikuti keterhubungan.
\end{proof}

Jika $S^n$ diberi orientasi, kelas teratas dapat dinormalisasi oleh
$\int_{S^n}\omega=1$. Teorema Stokes memastikan bahwa integral hanya
bergantung pada kelas kohomologi.

\subsection{Teorema de Rham dan gerbang topologi diferensial}

\begin{Theorem}[Teorema de Rham — pernyataan]\label{o011-br-thm-de-rham}
Integrasi bentuk pada rantai singular mulus menginduksi isomorfisme alami
\[
 H^k_{\mathrm{dR}}(M)\xrightarrow{\ \cong\ }
 H^k_{\mathrm{sing}}(M;\mathbb R)
\]
untuk setiap manifold mulus $M$ dan setiap $k$.
\end{Theorem}

Teorema ini dinyatakan, bukan dibuktikan, karena pembuktian lengkap memerlukan
kompleks singular, subdivisi, dan argumen berkas atau resolusi halus yang berada
di luar jembatan ini. Isi konseptualnya ialah bahwa perhitungan dengan bentuk
diferensial menghasilkan invariant topologis yang sama dengan kohomologi
singular berkoefisien nyata.

\paragraph{Derajat.}
Jika $f\colon M\to N$ mulus antara manifold tertutup, terhubung, berorientasi,
dan berdimensi sama $n$, tindakan $f^*$ pada kohomologi teratas adalah
perkalian oleh suatu bilangan bulat $\deg f$. Secara ekuivalen,
\[
 \int_M f^*\omega=(\deg f)\int_N\omega
\]
untuk setiap bentuk volume tertutup $\omega$. Dengan memilih nilai regular,
derajat dapat dihitung sebagai jumlah bertanda titik-titik prapeta.

\paragraph{Transversalitas.}
Jika $f\colon M\to N$ transversal terhadap submanifold $Z\subseteq N$, maka
$f^{-1}(Z)$ adalah submanifold berdimensi
$\dim M-\operatorname{codim}Z$. Transversalitas memungkinkan kelas kohomologi,
bilangan perpotongan, dan derajat direpresentasikan oleh geometri prapeta yang
stabil di bawah perturbasi kecil. Modul ini hanya memakai pernyataan gerbang
tersebut; teorema aproksimasi transversal memerlukan pembahasan tersendiri.

\paragraph{Teori Morse.}
Fungsi Morse mempunyai titik kritis nondegenerat. Ketika nilai melewati titik
kritis berindeks $\lambda$, sublevel berubah melalui pemasangan pegangan
$\lambda$. Dari sini lahir kompleks Morse, ketaksamaan Morse, dan hubungan
antara jumlah titik kritis dengan kohomologi. Pembuktian memerlukan lema Morse,
aliran gradien, dan analisis kekompakan; tidak satu pun disamarkan sebagai
telah dibuktikan di sini.

\begin{bemerkung}\label{o011-br-rem-roadmap}
Jembatan ini menutup rantai logis yang diperlukan untuk kurikulum dasar:
bentuk dan $d$ $\to$ kohomologi $\to$ invariansi homotopi $\to$
Mayer--Vietoris dan contoh $\to$ teorema de Rham sebagai identifikasi global
$\to$ gerbang derajat, transversalitas, serta Morse. Teori lanjutan dapat
memulai tepat dari tiga gerbang terakhir tanpa mengulang bagian sebelumnya.
\end{bemerkung}
